\section{Mon (Monoides Finitos)}

Para que uma construção em \textbf{Mon} seja válida, o objeto resultante deve manter a integridade da operação binária (associatividade) e garantir a existência de uma identidade global.

\subsection{Fundamentação Teórica}
\begin{itemize}
    \item \textbf{Produtos:} O produto de monoides define a operação componente a componente: $(m,n) * (m',n') = (m * m', n * n')$.
    \item \textbf{Equalizadores:} Corresponde ao sub-monoide que contém todos os elementos onde dois homomorfismos produzem o mesmo resultado.
\end{itemize}

\subsection{Implementação Computacional}
A construção do produto exige a geração de uma nova Tabela de Cayley expandida.

\begin{lstlisting}[caption={Criação da Tabela de Cayley para Produto de Monoides}]
function MonP = monoid_product(MA, idA, MB, idB)
    na = size(MA,1); nb = size(MB,1);
    MonP.table = zeros(na*nb, na*nb);
    
    % Construção da tabela via mapeamento de índices lineares
    for i=1:na*nb
        for j=1:na*nb
            [ia, ib] = ind2sub([na, nb], i);
            [ja, jb] = ind2sub([na, nb], j);
            % Operação componente a componente
            resA = MA(ia, ja);
            resB = MB(ib, jb);
            MonP.table(i,j) = sub2ind([na, nb], resA, resB);
        end
    end
    MonP.identity = sub2ind([na, nb], idA, idB);
end
\end{lstlisting}